\documentclass[12pt,reqno]{article}
\usepackage{amsmath,amssymb,amsfonts,amsthm}
\usepackage{booktabs}
\usepackage[hidelinks]{hyperref}

\theoremstyle{plain}
\newtheorem{theorem}{Theorem}
\newtheorem{lemma}[theorem]{Lemma}
\newtheorem{corollary}[theorem]{Corollary}
\newtheorem{proposition}[theorem]{Proposition}
\newtheorem{conjecture}[theorem]{Conjecture}
\newtheorem{openproblem}[theorem]{Open Problem}

\theoremstyle{definition}
\newtheorem{definition}[theorem]{Definition}
\newtheorem{example}[theorem]{Example}

\theoremstyle{remark}
\newtheorem{remark}[theorem]{Remark}


\newcommand{\N}{\mathbb{N}}
\newcommand{\Z}{\mathbb{Z}}
\newcommand{\np}{\operatorname{nextprime}}
\newcommand{\prevprime}{\operatorname{prevprime}}
\newcommand{\seqnum}[1]{\href{https://oeis.org/#1}{\rm \underline{#1}}}

\title{Computational Extension of OEIS Sequence A052075:
A New Term and an Exhaustive Search to $4 \times 10^{14}$}

\author{
Carlo Corti\\
Independent researcher\\
21-040 Kalin\'owka\\
Poland\\
\texttt{carlo.corti@outlook.com}
}

\date{}

\begin{document}

\maketitle

\begin{abstract}
We extend OEIS sequence \seqnum{A052075}, which consists of primes whose
successor prime appears as a contiguous decimal substring of the cube.
We carry out an exhaustive computation over the interval from
$6.9 \times 10^{12}$ to $4 \times 10^{14}$ using a segmented sieve of
Eratosthenes, arbitrary-precision arithmetic, and an exact modular
filter for trailing matches.
The search examines approximately $1.2 \times 10^{13}$ prime candidates
in 37 contiguous segments and finds exactly one new term,
$a(16) = 287\,902\,832\,031\,253$, the first addition since 2018.
This computation yields the improved lower bound
$a(17) > 4 \times 10^{14}$.
We also prove a congruential restriction on trailing matches modulo~10
and discuss a heuristic model suggesting that the sequence is infinite.
\end{abstract}

2020 Mathematics Subject Classification: Primary 11A41;
Secondary 11Y55, 11N05, 11A07.

Key words and phrases: prime number, integer sequence,
substring, cube of a prime, segmented sieve, computational number
theory.

\section{Introduction}
\label{sec:intro}

A prime $p$ belongs to sequence \seqnum{A052075} if the decimal
representation of its successor prime $q = \np(p)$ appears as a
contiguous substring of the decimal representation of $p^3$.
De Geest \cite{DeGeest2000} introduced the sequence in January~2000,
and Handler \cite{Handler2006} extended it to~15 terms in~2006.
The most recent OEIS entry \cite{OEIS_A052075} records the lower bound
$a(16) > 6.9 \times 10^{12}$ due to Resta \cite{Resta2018},
established in July~2018.

The computational difficulty has two distinct sources.
First, the number of primes below $X$ grows as $X / \ln X$ by the
prime number theorem, so an exhaustive search to $X$ requires testing
on the order of $X / \ln X$ candidates.
Second, computing $p^3$ for a 15-digit prime requires approximately
45-digit arithmetic, which overflows all standard fixed-width integer
types.
Specifically, for $p > 6.98 \times 10^{12}$ the cube exceeds the
capacity of \texttt{\_\_uint128\_t}; the exact threshold is
$(2^{128} - 1)^{1/3} \approx 6.981 \times 10^{12}$.
Therefore arbitrary-precision arithmetic is required throughout
the searched interval.

The present work resolves both difficulties.
We use the GNU Multiple Precision Arithmetic Library
(GMP) \cite{GMP} for exact 45-digit cube computation, and we introduce
a modular filter that detects trailing-position matches with very high
efficiency.
The resulting implementation exhausts $[6.9 \times 10^{12},
4 \times 10^{14}]$ in about 66 elapsed hours on a commodity workstation.

The main results are as follows.
\begin{enumerate}
\item $a(16) = 287\,902\,832\,031\,253$ (Theorem~\ref{thm:a16}),
      the first new term of \seqnum{A052075} since 2018.
\item $a(17) > 4 \times 10^{14}$ (Theorem~\ref{thm:a17}),
      extending Resta's bound by a factor of approximately~58.
\item A congruential classification of trailing matches in
      Proposition~\ref{prop:mod10}.
\item A heuristic discussion suggesting that \seqnum{A052075} is infinite.
\end{enumerate}

\subsection{Notation}
\label{sec:notation}

Throughout the paper, $p$ denotes a prime, $q = \np(p)$ its successor
prime, and $g = q - p$ the prime gap.
We write $\lfloor x \rfloor$ for the floor function and $\pi(x)$ for
the prime-counting function.
All logarithms are natural unless otherwise stated.
We use the notation $[L, R)$ for half-open integer intervals.
For any positive integer $n$, we denote by $s(n)$ the decimal string
of $n$ without leading zeros, and by
$d(n) = \lfloor \log_{10} n \rfloor + 1$ its digit count.

\section{Algorithm}
\label{sec:algorithm}

\subsection{Segmented sieve}
\label{sec:sieve}

Candidate primes are enumerated by a segmented sieve of
Eratosthenes \cite{Sorenson1991}.
The search interval $[L, R)$ is divided into macro-blocks of width
$W = 1.28 \times 10^{8}$.
Within each block, a standard byte-array sieve is applied after
precomputing all primes up to $\sqrt{R}$ by a preparatory sieve.
OpenMP \cite{OpenMP} parallelism is applied across macro-blocks.
With 16 logical threads on the test hardware, the implementation
achieves a throughput of $45$--$59 \times 10^6$ prime evaluations
per second, or approximately $3.1 \times 10^6$ per thread, decreasing
slowly with the digit count of $p$ because of the growing cost of
GMP arithmetic.

For primes $p \leq 6.98 \times 10^{12}$, the cube $p^3$ fits in
\texttt{\_\_uint128\_t} and no GMP arithmetic is needed.
For larger primes the implementation uses \texttt{mpz\_t} arithmetic.
Primality of candidates is guaranteed by the sieve construction;
independent verification of the new term uses deterministic
Miller--Rabin \cite{CrandallPomerance2005}.

\subsection{Substring matching}
\label{sec:matching}

Given a prime $p$ and $q = \np(p)$, a match occurs when the decimal
string $s(q)$ is a contiguous substring of $s(p^3)$.
Let $d_p = d(p)$ and $d_q = d(q)$ denote the digit counts of $p$
and $q$, respectively.
In the vast majority of cases in the searched interval, $d_q = d_p$.
The exception occurs for the finitely many primes whose successor
prime crosses a power-of-10 boundary (e.g., $\np(9\,999\,999\,999\,971)
= 10\,000\,000\,000\,037$, so $d_q = d_p + 1$).
The implementation uses the exact value of $d_q$ in all cases.
The digit count of $p^3$ satisfies
$d(p^3) \in \{3d_p - 2,\; 3d_p - 1,\; 3d_p\}$;
we write $D = d(p^3)$ for the exact value.

We call the match \emph{trailing} if $s(q)$ occupies the last $d_q$
digits of $s(p^3)$, that is, if $p^3 \equiv q \pmod{10^{d_q}}$.
We call it \emph{leading} if $s(q)$ occupies the first $d_q$ digits,
and \emph{interior} otherwise.
Among the 16 known terms, trailing matches occur for $n = 10, 12, 15,
16$, and leading matches occur for $n = 1, 2, 13$;
see Table~\ref{tab:oracle}.

\subsection{Modular filter for trailing matches}
\label{sec:filter}

The trailing condition $p^3 \equiv q \pmod{10^{d_q}}$ can be rewritten
purely in terms of $p$ and the prime gap $g = q - p$ as
\begin{equation}
 p^3 \equiv p + g \pmod{10^{d_q}}.
\label{eq:trailing}
\end{equation}
This simplifies to
\begin{equation}
 p(p^2 - 1) \equiv g \pmod{10^{d_q}}.
\label{eq:filter}
\end{equation}

\begin{lemma}
\label{lem:filter}
A prime $p$ with $\np(p) = p + g$ satisfies the trailing-match
condition if and only if $p(p^2 - 1) \equiv g \pmod{10^{d_q}}$,
where $d_q = \lfloor \log_{10}(p + g) \rfloor + 1$.
\end{lemma}

\begin{proof}
The trailing condition $p^3 \equiv q \pmod{10^{d_q}}$ is equivalent
to $p^3 - q \equiv 0 \pmod{10^{d_q}}$.
Substituting $q = p + g$ gives
$p^3 - p - g \equiv 0 \pmod{10^{d_q}}$,
which is equivalent to $p(p^2 - 1) \equiv g \pmod{10^{d_q}}$.
\end{proof}

The filter of Lemma~\ref{lem:filter} yields an immediate structural
consequence.

\begin{proposition}
\label{prop:mod10}
Let $p > 5$ be a prime with $\np(p) = p + g$.
If $p$ admits a trailing match, then
\[
 g \equiv \begin{cases}
 0 \pmod{10}, & \text{if } p \equiv 1 \pmod{10}, \\
 4 \pmod{10}, & \text{if } p \equiv 3 \pmod{10}, \\
 6 \pmod{10}, & \text{if } p \equiv 7 \pmod{10}, \\
 0 \pmod{10}, & \text{if } p \equiv 9 \pmod{10}.
 \end{cases}
\]
\end{proposition}

\begin{proof}
By Lemma~\ref{lem:filter}, a trailing match requires
$p(p^2 - 1) \equiv g \pmod{10^{d_q}}$, and hence
$p(p^2 - 1) \equiv g \pmod{10}$.
Since $p > 5$ is prime, $p \bmod 10 \in \{1, 3, 7, 9\}$.
We compute $p(p^2 - 1) \bmod 10$ for each residue class:
\begin{align*}
 p \equiv 1 \pmod{10} &\implies p(p^2 - 1) \equiv 1 \cdot 0 = 0 \pmod{10}, \\
 p \equiv 3 \pmod{10} &\implies p(p^2 - 1) \equiv 3 \cdot 8 = 24 \equiv 4 \pmod{10}, \\
 p \equiv 7 \pmod{10} &\implies p(p^2 - 1) \equiv 7 \cdot 48 = 336 \equiv 6 \pmod{10}, \\
 p \equiv 9 \pmod{10} &\implies p(p^2 - 1) \equiv 9 \cdot 80 = 720 \equiv 0 \pmod{10}.
\end{align*}
This proves the claim.
\end{proof}

\begin{remark}
Proposition~\ref{prop:mod10} is consistent with all four trailing
matches among the 16 known terms.
For $a(10)$ we have $p \equiv 7$ and $g = 56 \equiv 6 \pmod{10}$;
for $a(12)$ we have $p \equiv 7$ and $g = 96 \equiv 6 \pmod{10}$;
for $a(15)$ we have $p \equiv 9$ and $g = 40 \equiv 0 \pmod{10}$;
and for $a(16)$ we have $p \equiv 3$ and $g = 24 \equiv 4 \pmod{10}$.
\end{remark}

\begin{remark}
The classification can be extended to higher moduli by applying the
Chinese remainder theorem to the factors $2^k$ and $5^k$ of $10^k$.
For example, modulo $100$ one obtains a finer partition of admissible
residue pairs $(p \bmod 100, g \bmod 100)$, further restricting the
candidates that can produce a trailing match.
\end{remark}

The modular filter of Lemma~\ref{lem:filter} is exact for trailing
matches: it accepts a candidate if and only if a trailing match
occurs.
In the implementation, the filter is evaluated as a fast 128-bit
modular computation for every prime $p$.
Candidates that pass the filter are flagged as potential trailing
matches and recorded immediately, without waiting for the full
substring scan.
Candidates that do not pass the filter are not discarded: every prime
in the search interval proceeds to the full computation of $p^3$ and
the complete scan of $s(p^3)$ for $s(q)$, so leading and interior
matches are not missed.

As a rough heuristic sanity check, for 15-digit primes the average
prime gap near $10^{15}$ is approximately $\ln(10^{15}) \approx 35$
by the prime number theorem.
This gives a theoretical expected number of trailing-match candidates
of order $(35 / 10^{15}) \times 1.2 \times 10^{13} \approx 0.4$
over the searched interval, consistent with finding exactly one such
candidate ($a(16)$) and zero spurious trailing hits.

\section{Computational results}
\label{sec:results}

\subsection{Oracle verification}
\label{sec:oracle}

All 15 previously known terms were verified independently at the start
of every search segment.
For each term $a(n)$, we confirmed the primality of $p$, the
correctness of $\np(p)$, the exact value of $p^3$, and the substring
match at the stated position.
Table~\ref{tab:oracle} lists the complete oracle set, including the
new term $a(16)$.

\begin{table}[ht]
\begin{center}
\small
\begin{tabular}{rllrrl}
\toprule
$n$ & $a(n) = p$ & $\np(p)$ & pos & len & type \\
\midrule
1  & 11                & 13                & 0  & 4  & L \\
2  & 101               & 103               & 0  & 7  & L \\
3  & 2\,239            & 2\,243            & 2  & 11 & I \\
4  & 34\,297           & 34\,301           & 2  & 14 & I \\
5  & 43\,789           & 43\,793           & 4  & 14 & I \\
6  & 53\,549           & 53\,551           & 1  & 15 & I \\
7  & 535\,487          & 535\,489          & 1  & 18 & I \\
8  & 59\,897\,017      & 59\,897\,039      & 14 & 24 & I \\
9  & 430\,784\,719     & 430\,784\,737     & 3  & 26 & I \\
10 & 2\,549\,592\,677  & 2\,549\,592\,733  & 19 & 29 & T \\
11 & 2\,837\,138\,669  & 2\,837\,138\,677  & 1  & 29 & I \\
12 & 97\,969\,345\,967 & 97\,969\,346\,063 & 22 & 33 & T \\
13 & 100\,000\,000\,019 & 100\,000\,000\,057 & 0 & 34 & L \\
14 & 328\,096\,840\,219 & 328\,096\,840\,223 & 20 & 35 & I \\
15 & 4\,110\,739\,763\,869 & 4\,110\,739\,763\,909 & 25 & 38 & T \\
\midrule
16 & 287\,902\,832\,031\,253 & 287\,902\,832\,031\,277 & 29 & 44 & T \\
\bottomrule
\end{tabular}
\caption{Complete known terms of \seqnum{A052075}.
Column pos gives the 0-based starting position of $\np(p)$ in $s(p^3)$;
column len gives the digit count $D = d(p^3)$;
column type classifies the match as leading (L), trailing (T),
or interior (I).
Full values of $p^3$ for all 16 terms are available in the oracle CSV
file at the GitHub repository.}
\label{tab:oracle}
\end{center}
\end{table}

\subsection{New term: \texorpdfstring{$a(16)$}{a(16)}}
\label{sec:a16}

The exhaustive search of the interval $[6.9 \times 10^{12},
4 \times 10^{14}]$ yields exactly one new term.

\begin{theorem}
\label{thm:a16}
We have $a(16) = 287\,902\,832\,031\,253$.
\end{theorem}

\begin{proof}
We verify the four defining properties.
\begin{enumerate}
\item $p = 287\,902\,832\,031\,253$ is prime, as confirmed by
      \texttt{sympy.isprime}.
\item $q = \np(p) = 287\,902\,832\,031\,277$, with gap $g = 24$,
      as confirmed by \texttt{sympy.nextprime}.
      This also shows that every integer in the open interval
      $(p, q)$ is composite.
\item The cube is the 44-digit integer
      \[
      p^3 = 23\,863\,701\,656\,637\,949\,988\,435\,953\,863\,
      287\,902\,832\,031\,277,
      \]
      as confirmed by GMP and by Python \texttt{int} arithmetic.
\item The decimal string of $q$ occupies positions~$29$ through~$43$
      of $s(p^3)$, using 0-based indexing.
      Thus the last 15 digits of the cube equal $q$, so this is a
      trailing match.
\end{enumerate}
The campaign summary in Tables~\ref{tab:campaign_a}
and~\ref{tab:campaign_b} documents the exhaustiveness of the search
on $[6.9 \times 10^{12}, 4 \times 10^{14}]$.
\end{proof}

\begin{remark}
The modular filter in \eqref{eq:filter} confirms
$p(p^2 - 1) \equiv 24 \pmod{10^{15}}$.
Indeed, $p \equiv 287\,902\,832\,031\,253 \pmod{10^{15}}$,
so $p(p^2 - 1) \bmod 10^{15} = 24$.
\end{remark}

\subsection{Improved lower bound}
\label{sec:lowerbound}

\begin{theorem}
\label{thm:a17}
We have $a(17) > 4 \times 10^{14}$.
\end{theorem}

\begin{proof}
This is a computational theorem: the conclusion rests on the verified
exhaustive search documented in Tables~\ref{tab:campaign_a}
and~\ref{tab:campaign_b}.
That search covers every prime in $[6.9 \times 10^{12},
4 \times 10^{14}]$ and detects no match of any type---leading,
trailing, or interior---other than the single term $a(16)$ found in
segment~S28.
By Theorem~\ref{thm:a16}, this term equals
$287\,902\,832\,031\,253$, which lies in
$[6.9 \times 10^{12}, 4 \times 10^{14}]$.
Therefore $a(17)$, the next term after $a(16)$, must exceed
$4 \times 10^{14}$.
In particular, segments S29 through S37 cover
$[2.92 \times 10^{14}, 4 \times 10^{14}]$ with zero hits,
confirming the absence of any additional term in that sub-interval.
\end{proof}

\subsection{Campaign summary}
\label{sec:campaign}

The search was executed as 37 contiguous segments on a workstation
with an AMD Ryzen~9 7940HS processor (8~cores, 16~threads),
64~GB DDR5 RAM, running Linux Mint~22.3.
Tables~\ref{tab:campaign_a} and~\ref{tab:campaign_b} summarize the
key parameters of each segment.

\begin{table}[ht]
\begin{center}
\small
\begin{tabular}{rlrrrl}
\toprule
Seg & Range $[L, R) \times 10^{12}$ & Primes ($\times 10^9$)
    & Elapsed (s) & Hits & Date \\
\midrule
S01 & $[6.9,\;14.9)$    & 266.7 & 4522 & 0 & 2026-03-24 \\
S02 & $[14.9,\;22.9)$   & 261.8 & 4576 & 0 & 2026-03-24 \\
S03 & $[22.9,\;30.9)$   & 258.7 & 4624 & 0 & 2026-03-24 \\
S04 & $[30.9,\;38.9)$   & 256.6 & 4688 & 0 & 2026-03-25 \\
S05 & $[38.9,\;46.9)$   & 254.9 & 4687 & 0 & 2026-03-25 \\
S06 & $[46.9,\;54.9)$   & 253.5 & 4647 & 0 & 2026-03-25 \\
S07 & $[54.9,\;62.9)$   & 252.3 & 4695 & 0 & 2026-03-25 \\
S08 & $[62.9,\;70.9)$   & 251.3 & 4633 & 0 & 2026-03-25 \\
S09 & $[70.9,\;78.9)$   & 250.4 & 4648 & 0 & 2026-03-25 \\
S10 & $[78.9,\;86.9)$   & 249.6 & 4714 & 0 & 2026-03-25 \\
S11 & $[86.9,\;94.9)$   & 248.9 & 4737 & 0 & 2026-03-26 \\
S12 & $[94.9,\;100.0)$  & 158.3 & 3051 & 0 & 2026-03-26 \\
S13 & $[100.0,\;112.0)$ & 371.6 & 7030 & 0 & 2026-03-26 \\
S14 & $[112.0,\;124.0)$ & 370.4 & 7198 & 0 & 2026-03-26 \\
S15 & $[124.0,\;136.0)$ & 369.3 & 7074 & 0 & 2026-03-26 \\
S16 & $[136.0,\;148.0)$ & 368.3 & 7172 & 0 & 2026-03-27 \\
S17 & $[148.0,\;160.0)$ & 367.3 & 7152 & 0 & 2026-03-27 \\
S18 & $[160.0,\;172.0)$ & 366.5 & 7432 & 0 & 2026-03-27 \\
S19 & $[172.0,\;184.0)$ & 365.7 & 7141 & 0 & 2026-03-27 \\
\bottomrule
\end{tabular}
\caption{Campaign summary, segments S01--S19.
Segment widths are $8 \times 10^{12}$ for S01--S11
and $12 \times 10^{12}$ for S13--S19
(S12 has width $5.1 \times 10^{12}$).
The time column reports elapsed wall-clock time of the C~process
with 16 OpenMP threads.
Prime counts are rounded to one decimal place; the total in
Table~\ref{tab:campaign_b} is computed from unrounded raw counts.}
\label{tab:campaign_a}
\end{center}
\end{table}

\begin{table}[ht]
\begin{center}
\small
\begin{tabular}{rlrrrl}
\toprule
Seg & Range $[L, R) \times 10^{12}$ & Primes ($\times 10^9$)
    & Elapsed (s) & Hits & Date \\
\midrule
S20 & $[184.0,\;196.0)$ & 365.0 & 7302 & 0 & 2026-03-28 \\
S21 & $[196.0,\;208.0)$ & 364.3 & 7144 & 0 & 2026-03-28 \\
S22 & $[208.0,\;220.0)$ & 363.7 & 7290 & 0 & 2026-03-28 \\
S23 & $[220.0,\;232.0)$ & 363.1 & 7343 & 0 & 2026-03-28 \\
S24 & $[232.0,\;244.0)$ & 362.5 & 7471 & 0 & 2026-03-29 \\
S25 & $[244.0,\;256.0)$ & 362.0 & 7377 & 0 & 2026-03-29 \\
S26 & $[256.0,\;268.0)$ & 361.5 & 7387 & 0 & 2026-03-29 \\
S27 & $[268.0,\;280.0)$ & 361.0 & 7406 & 0 & 2026-03-29 \\
S28 & $[280.0,\;292.0)$ & 360.5 & 7480 & 1 & 2026-03-30 \\
S29 & $[292.0,\;304.0)$ & 360.1 & 7412 & 0 & 2026-03-30 \\
S30 & $[304.0,\;316.0)$ & 359.6 & 7970 & 0 & 2026-03-30 \\
S31 & $[316.0,\;328.0)$ & 359.2 & 7352 & 0 & 2026-03-30 \\
S32 & $[328.0,\;340.0)$ & 358.8 & 7345 & 0 & 2026-03-30 \\
S33 & $[340.0,\;352.0)$ & 358.5 & 7434 & 0 & 2026-03-31 \\
S34 & $[352.0,\;364.0)$ & 358.1 & 7493 & 0 & 2026-03-31 \\
S35 & $[364.0,\;376.0)$ & 357.7 & 7361 & 0 & 2026-03-31 \\
S36 & $[376.0,\;388.0)$ & 357.4 & 7434 & 0 & 2026-03-31 \\
S37 & $[388.0,\;400.0)$ & 357.1 & 7446 & 0 & 2026-03-31 \\
\midrule
Total & $[6.9,\;400.0)$ & $12\,031.9$ & $237\,868$ & 1 & \\
\bottomrule
\end{tabular}
\caption{Campaign summary, segments S20--S37, and totals.
Elapsed times in seconds; the total row gives the sum over all 37
segments.
Prime counts are rounded to one decimal place; the total is computed
from unrounded raw counts.}
\label{tab:campaign_b}
\end{center}
\end{table}

The total elapsed wall-clock time is $237\,868$ seconds,
or about 66.1 hours.
With 16 OpenMP threads, the average throughput is approximately
$50 \times 10^6$ prime evaluations per second, corresponding to
roughly $3.1 \times 10^6$ evaluations per thread.
The calendar elapsed time is approximately 7.2 days, from
2026-03-24 to 2026-03-31, reflecting discontinuous use of the
workstation.

\subsection{Independent verification}
\label{sec:verification}

The new term $a(16) = 287\,902\,832\,031\,253$ was verified
independently in Python~3 with \texttt{sympy}
(version~1.14.0) as follows.

\begin{enumerate}
\item \texttt{sympy.isprime(287902832031253)} returns \texttt{True}.
\item \texttt{sympy.nextprime(287902832031253)} returns
      \texttt{287902832031277}.
\item \texttt{287902832031253**3} evaluates to\\
      \texttt{23863701656637949988435953863287902832031277}.
\item \texttt{str(287902832031277) in str(287902832031253**3)}
      returns \texttt{True}, at position~29 of the 44-digit string.
\end{enumerate}

All four conditions are satisfied.
For numbers below $2^{64} \approx 1.84 \times 10^{19}$,
the \texttt{sympy.isprime} function performs a deterministic
primality test based on trial division and deterministic
Miller--Rabin with verified witness sets \cite{SymPy},
so the verification in step~1 is conclusive rather than probabilistic.
The value $p \approx 2.88 \times 10^{14}$ lies well below this bound.

\section{Probabilistic heuristic}
\label{sec:heuristic}

\subsection{Expected density of terms}
\label{sec:density}

For a prime $p$ with $d = d(p)$ decimal digits, let $q = \np(p)$,
which also has $d$ digits for all sufficiently large $p$ in a fixed
decade.
The cube $p^3$ has $D \in \{3d - 2,\; 3d - 1,\; 3d\}$ digits.
For a rough estimate we take $D \approx 3d$.

The probability that a uniformly chosen $d$-digit string $s(q)$
appears as a substring of a uniformly chosen $D$-digit string
$s(p^3)$ is approximately
\begin{equation}
P_{\mathrm{sub}} \approx \frac{D - d + 1}{10^d} \approx \frac{2d}{10^d},
\label{eq:psub}
\end{equation}
since there are approximately $D - d + 1 \approx 2d$ possible starting
positions and each position matches with probability $10^{-d}$.

The number of $d$-digit primes is approximately
$\pi(10^d) - \pi(10^{d-1}) \approx 8 \cdot 10^{d-1} / d$
(a standard rough estimate from the prime number theorem).
The expected number of terms of \seqnum{A052075} in the $d$-digit
decade is therefore
\begin{equation}
E_d \approx \frac{8 \cdot 10^{d-1}}{d} \cdot \frac{2d}{10^d} = 1.6.
\label{eq:Ed_raw}
\end{equation}

The strings $s(q)$ and $s(p^3)$ are not independent, because both are
determined by $p$.
The study of digit strings in prime-related sequences \cite{Dudek2014}
suggests that such arithmetic correlations can produce a moderate
downward bias.
Empirically, the 15 known terms in the 12 decades from $d = 2$ through
$d = 13$ give an observed density of $15/12 = 1.25$ terms per decade.
The ratio of this observed density to the crude estimate
$E_d \approx 1.6$ yields a correction factor of
\[
\frac{15}{12 \times 1.6} \approx 0.78.
\]
Applying this correction factor to the crude estimate gives an
adjusted expected count of $0.78 \times 1.6 = 1.25$ terms per
decade, consistent with the observed density.

\subsection{Implication for infinitude}
\label{sec:infinitude}

Since the adjusted estimate of approximately 1.25 terms per decade
remains bounded away from zero across decades, the heuristic model
suggests that sufficiently large decades should continue to contain
terms of \seqnum{A052075} with positive probability.
Under an additional assumption that occurrences in well-separated
decades are approximately independent---which is plausible but
unproved---this points toward infinitely many terms.
This is only a heuristic argument.
The events are not independent across primes, and the
approximation~\eqref{eq:psub} treats $s(q)$ and $s(p^3)$ as
independent, which is false.
The discussion serves only as motivation for the following conjecture.

\begin{conjecture}
\label{conj:infinite}
The sequence \seqnum{A052075} is infinite.
\end{conjecture}

A rigorous proof of Conjecture~\ref{conj:infinite} seems to be beyond
current methods, since the problem mixes the distribution of prime gaps
with the non-algebraic structure of decimal representations.

\section{Conclusion and open problems}
\label{sec:conclusion}

We determine $a(16) = 287\,902\,832\,031\,253$, the first new term of
\seqnum{A052075} since 2018, and establish the improved lower bound
$a(17) > 4 \times 10^{14}$.
The computation requires about 66 elapsed hours on a commodity
workstation and is made feasible by GMP arithmetic, an exact modular
filter for trailing-position candidates, and an OpenMP-parallelized
segmented sieve.

The main open problems are as follows.

\begin{openproblem}
\label{op:a17}
Determine $a(17)$.
Under the heuristic model, the interval $[4 \times 10^{14}, 10^{15}]$
has logarithmic length $\log_{10}(2.5) \approx 0.398$ decades.
Applying the observed density of 1.25 terms per decade gives an
expected count of about $0.398 \times 1.25 \approx 0.50$, so the
Poisson-model probability of at least one term in this interval
is approximately $1 - e^{-0.50} \approx 0.39$.
\end{openproblem}

\begin{openproblem}
\label{op:infinite}
Prove or disprove Conjecture~\ref{conj:infinite}.
\end{openproblem}

\begin{openproblem}
\label{op:positions}
Characterize which positions, namely leading, trailing, and interior,
can support a match, and determine whether trailing matches are denser
than interior matches.
Among the 16 known terms, 4 are trailing, 3 are leading, and
9 are interior.
\end{openproblem}

\begin{openproblem}
\label{op:analogues}
Investigate the analogous sequences for $p^k$ with $k \geq 2$,
including the square analogue
\seqnum{A052073} \cite{OEIS_A052073},
the related sequence \seqnum{A321796} \cite{OEIS_A321796},
in which $\prevprime(p)$ is a substring of $p^3$,
and the sequence \seqnum{A381969} \cite{OEIS_A381969},
in which $\prevprime(p)$ is a substring of $p^2$.
\end{openproblem}

\section{Acknowledgments}

The author thanks the OEIS contributors whose earlier computations
established the initial terms and the previous lower bound for
\seqnum{A052075}.

\section{Data availability}

The complete dataset, including the oracle CSV file, per-segment logs,
campaign summary, and machine-readable results, is available at
the GitHub repository
\url{https://github.com/carcorti/OEIS-A052075}
and at the corresponding Zenodo archive (DOI to be updated before
final submission).
The results are also available through the OEIS entry
for \seqnum{A052075}.

\section{Code availability}

The C source file \texttt{a052075\_v4.c} (version~4 of the search
program), the associated Makefile, and a Python verification script
\texttt{verify\_a052075.py} are available at the GitHub and Zenodo
repositories listed in the previous section.

\section{AI use disclosure}

The author used large language model tools during the preparation of
this work.
The AI tools assisted with code review, verification of computational
results, and identifying weaknesses in the exposition.
No prose in this manuscript was generated by a large language model.
All mathematical arguments, computational claims, and final wording
of the manuscript were produced and verified by the author, who takes
full scientific responsibility for the paper.

\section{Sequences}

The sequences appearing in this paper are as follows.

\begin{itemize}
\item \seqnum{A052073}, the square analogue of \seqnum{A052075}
      mentioned in Open Problem~\ref{op:analogues}.
\item \seqnum{A052075}, primes $p$ such that the successor prime
      $\np(p)$ appears as a contiguous decimal substring of $p^3$.
\item \seqnum{A321796}, primes $p$ such that the predecessor prime
      $\prevprime(p)$ appears as a contiguous decimal substring
      of $p^3$, mentioned in Open Problem~\ref{op:analogues}.
\item \seqnum{A381969}, primes $p$ such that the predecessor prime
      $\prevprime(p)$ appears as a contiguous decimal substring
      of $p^2$, mentioned in Open Problem~\ref{op:analogues}.
\end{itemize}

\bibliographystyle{jis}
\bibliography{A052075_rev11}

\end{document}
